\chapter{Literature Review}

\section{Foundational Theory}

\subsection{Gig Economy and Non-Traditional Workers}

The gig economy is a digital platform-based economic model that facilitates temporary or task-based work, where workers earn income from short-term projects without permanent employment ties \cite{wood2019goodgig}. According to the International Labour Organization (ILO), the gig economy is characterized by high flexibility, worker autonomy, but also income uncertainty and a lack of social protections \cite{ilo2025worldemployment}. In Indonesia, the number of informal workers (including gig workers and part-timers) reached approximately 86.58 million people in February 2025 (59.4\% of the total workforce according to BPS), with subsets of platform-based gig work---such as freelancing in graphic design, website development, photography, and others---continuing to grow rapidly among young people \cite{bps2025tenagakerja}. The majority of these workers are young, and many are students who engage in side gigs or part-time jobs to supplement their income.

Non-traditional workers, particularly students, often face a double burden between academic responsibilities and work, leading to high psychological pressure \cite{syiahkuala2024burnout}. Key characteristics of gig work include the use of platform algorithms for task allocation, rating systems that influence earnings, and social isolation due to limited face-to-face interactions \cite{frontiers2025gigwork}. These conditions make gig workers vulnerable to mental health issues, such as burnout. This foundational theory serves as the basis for understanding the context of non-traditional workers in Indonesia, which will later be linked to burnout risks and potential digital interventions.

\subsection{Burnout}

Burnout is defined as a psychological syndrome characterized by emotional exhaustion, depersonalization (cynical attitudes toward work), and reduced personal accomplishment (a diminished sense of self-efficacy) \cite{maslach2016burnout}. The most commonly used framework for measuring burnout is the Maslach Burnout Inventory (MBI), developed by Christina Maslach, which assesses these three main dimensions \cite{maslach1996mbi}.

In the gig economy, burnout is often triggered by factors such as job insecurity, algorithmic management (pressure from ratings and platform algorithms), and a lack of work-life balance \cite{pmc2025scopingreview}. Studies indicate that gig workers face a 50--60\% higher risk of burnout compared to permanent employees, with symptoms including chronic fatigue and decreased motivation \cite{jmir2025digitalinterventions}. In Indonesia, students with side gigs or part-time jobs also experience high burnout rates, reaching up to 91\% in some local samples, due to schedule conflicts between studies and work, as well as algorithmic pressures from platforms \cite{syiahkuala2024burnout, rimtb2024burnout}.

\subsection{Mental Health and mHealth Applications}

Mental health refers to a state of emotional, psychological, and social well-being that enables individuals to cope with daily stressors \cite{who2022mentalhealth}. In the digital era, mobile health (mHealth) applications have become increasingly popular as self-help interventions, including chatbots and mood trackers for emotion journaling and early stress pattern detection \cite{jmir2025chatbot}. Examples such as Woebot and Wysa use chatbots to provide conversational support and mood tracking, which have been shown to reduce anxiety symptoms by up to 20\% among regular users \cite{fitzpatrick2017woebot}. Recent trends indicate increased mHealth adoption post-pandemic, particularly for stigma-free self-tracking, with high effectiveness among young populations.

Mood trackers are a core feature in mHealth apps, allowing users to log daily emotions, identify triggers, and view trends through simple visualizations (e.g., line graphs) \cite{torous2020moodtracking}. This feature is effective for burnout management as it supports self-awareness without requiring direct professional intervention.

\subsection{Basic Principles of UI/UX}

\subsubsection{Context-Aware and Adaptive Interfaces}

Adaptive UI/UX, or a context-aware interface, is an extension of standard usability principles that allows an interface to dynamically adjust itself based on user context, including available time, location, stress levels, physical activity, or input preferences \cite{benyon2014interactive}. According to Calvary et al. \cite{calvary2003plasticity}, this adaptation encompasses the plasticity of user interfaces, where the system can automatically alter its layout, content, or interaction modes to enhance flexibility and efficiency of use (aligning with Nielsen's sixth heuristic). In the context of non-traditional workers with irregular schedules and high cognitive loads, adaptations such as switching to one-tap emojis during busy periods, offering voice input while mobile (e.g., driving or multitasking in a side gig), and simplifying trend visualizations during high stress are crucial for reducing cognitive load and increasing emotional engagement. This approach remains grounded in User-Centered Design but adds a contextual layer that bridges the gaps found in conventional mHealth applications.

\subsubsection{Cognitive Load and Hick's Law in UI Design}

The necessity for an adaptive interface in this study is theoretically driven by John Sweller's Cognitive Load Theory, specifically the need to minimize extraneous cognitive load---the mental effort required to process poorly designed or overly complex interfaces \cite{sweller1988cogload}. For gig workers experiencing burnout, their working memory capacity is already severely depleted. To address this, the adaptive UI relies heavily on the Hick--Hyman Law (Hick's Law), which states that the time it takes for a person to make a decision increases logarithmically with the number and complexity of choices presented. By intentionally restricting the number of interactive elements, buttons, and navigation menus during high-stress periods, the application minimizes the decision-making time, thereby mathematically reducing the user's extraneous cognitive load and lowering the barrier to daily mood logging.

\subsubsection{Usability Principles and Iterative Design Framework}

Human-Computer Interaction (HCI) is a field of study that examines the interactions between humans and computer systems, aiming to create usable, useful, and user-friendly interfaces \cite{dix2004hci}. The usability principles proposed by Jakob Nielsen serve as a primary foundation in UI/UX design, consisting of ten heuristics \cite{nielsen1993usability}:

\begin{enumerate}
	\item Visibility of system status (users always know what is happening).
	\item Match between system and the real world (language and concepts align with the user's real world).
	\item User control and freedom (users have control and can undo actions).
	\item Consistency and standards (consistent with general conventions).
	\item Error prevention, recognition, and recovery.
	\item Flexibility and efficiency of use (suitable for both novices and experts).
	\item Aesthetic and minimalist design.
	\item Help users recognize, diagnose, and recover from errors.
	\item Help and documentation.
\end{enumerate}

To apply these heuristics effectively, this study adopts a Literature-Driven Iterative Design approach inspired by human-centered design principles. Traditionally, design cycles involve understanding the context of use, specifying user requirements, designing solutions, and evaluating against those requirements \cite{iso9241210}. In this study, the initial understanding of user context and requirements is achieved by synthesizing established academic literature regarding the cognitive load and burnout of non-traditional workers. This ensures the interface is theoretically grounded to support users experiencing high stress. Subsequently, the evaluation phase directly involves users in controlled settings to test the prototype's efficiency and emotional impact against those requirements.

These principles form the basis for designing the UI/UX of the mood tracker application, where good usability can enhance user engagement and the effectiveness of self-help interventions for burnout management. However, for users in specific contexts like non-traditional workers, these standard heuristics must be extended into adaptive UI/UX, as discussed in the technical concepts below.

\section{Technical Concepts}

\subsection{Mood Tracker in Mental Health Applications}

A mood tracker is a feature in mobile health (mHealth) applications that enables users to log their daily emotional states, identify stress triggers, and visualize emotional patterns through simple graphs \cite{torous2020moodtracking}. The primary functions of a mood tracker include:

\begin{enumerate}
	\item \textbf{Emotion logging:} Users select or write their mood condition (e.g., stressed, calm, tired) along with triggers (e.g., work deadlines or traffic jams).
	\item \textbf{Trend visualization:} Displaying data in the form of line graphs or pie charts to observe burnout patterns over specific periods.
	\item \textbf{Basic insights:} Providing simple recommendations based on patterns, such as ``You often feel stressed at night; try resting earlier.''
\end{enumerate}

In the context of non-traditional workers, mood trackers have significant potential as effective self-help tools, allowing burnout monitoring without direct professional intervention \cite{jmir2025digitalinterventions}. Applications such as Daylio or Moodpath have demonstrated that regular mood tracking can increase self-awareness by up to 25\% among users with chronic stress symptoms \cite{fitzpatrick2017woebot}. Integration with prebuilt chatbots (e.g., Gemma 3) enables the mood tracker to deliver simple conversational responses, such as ``Thank you for sharing; what was the trigger today?'' Adaptive examples include a simplified UI when the user is in busy hours.

\subsection{Usability Metrics for UI/UX Evaluation}

Evaluating UI/UX in mental health applications requires metrics that not only measure ease of use but also emotional impact and personalization. This study employs three adapted and specifically developed metrics for the context of non-traditional workers with high cognitive load.

\subsubsection{Time on Task (ToT)}

Time on Task (ToT) is the standard usability metric for measuring interaction efficiency, replacing completion rate metrics when the primary goal is assessing cognitive load reduction. Based on the ISO 9241-11 usability framework, efficiency is quantified by the exact duration (in seconds) a user requires to successfully complete a specific scenario, such as logging a mood and its trigger \cite{iso92411}.

By utilizing a within-subjects A/B testing design, this study compares the ToT of the standard interface against the adaptive ``Rush Hour'' interface. A statistically significant reduction in ToT serves as empirical evidence that the adaptive UI successfully minimizes extraneous cognitive load for non-traditional workers.

\subsubsection{Emotional Response Index (ERI)}

ERI assesses the user's affective response to the interface, a crucial component for mental health applications designed for high-stress demographics. To ensure construct validity, the ERI questionnaire items are theoretically grounded in the emotional valence dimensions of the mHealth App Usability Questionnaire (MAUQ) \cite{zhou2019mauq} and the Positive and Negative Affect Schedule (PANAS) \cite{watson1988panas}. Evaluated on a 5-point Likert scale, items measure subjective feelings of relief, support, and reduced anxiety post-interaction, ensuring usability evaluations extend beyond mechanical efficiency to include psychological impact. The ERI items are as follows:

\begin{enumerate}
	\item The application's interface helps me process my emotions without feeling mentally drained.
	\item The simplified visual design of the application makes me feel calm even when I am in a rush.
	\item The chatbot's empathetic responses made me feel heard and supported.
	\item I felt a sense of relief after logging my mood using the ``Rush Hour'' feature.
	\item The application provides a safe and comfortable space for me to reflect on my daily stress triggers.
	\item Interacting with the application during a busy shift did not increase my anxiety or frustration.
	\item The chatbot's conversational tone is soothing and appropriate for someone experiencing burnout.
	\item Overall, interacting with this application positively impacts my emotional well-being.
\end{enumerate}

\subsubsection{Personalization Effectiveness Index (PEI)}

PEI evaluates the system's ability to dynamically align with the user's contextual needs. To ensure reliability, the evaluation items for this metric are adapted from the Perceived Personalization Scale commonly used in HCI research \cite{komiak2006personalization}. Assessed via a 5-point Likert scale, the PEI measures the user's perception of the system's contextual awareness, evaluating factors such as the relevance of the Gemma 3 LLM prompts and the timeliness of the adaptive UI triggers (e.g., anticipating the preferred input method during busy periods). The PEI items are as follows:

\begin{enumerate}
	\item The ``Rush Hour'' interface accurately anticipates my need for a faster input method during my busy shifts.
	\item The chatbot's responses felt specifically tailored to my unique mood history rather than being generic or automated.
	\item The application effectively accommodates my irregular schedule as a student with a side gig.
	\item The adaptive input features (e.g., voice logging, one-tap emojis) are highly relevant to my physical context when I am working.
	\item The chatbot provides conversational insights that are directly applicable to my personal burnout triggers.
	\item The application feels like it understands my personal time limitations and cognitive exhaustion.
	\item During high-stress periods, the interface successfully hides unnecessary features and only shows me what I need.
	\item The overall experience with the application feels highly personalized to my individual working conditions.
\end{enumerate}

\subsection{Adaptive UI/UX as an Extension of Standard Principles}

Although Nielsen's usability heuristics and foundational design principles provide a strong baseline, users with high cognitive demands---such as non-traditional workers---require the development of adaptive UI/UX. Adaptive UI dynamically adjusts the interface based on user context, such as available time, stress levels, or input preferences \cite{pmc2025scopingreview}. Relevant adaptation examples for mood trackers include:

\begin{itemize}
	\item \textbf{Adaptive user interface:} Adjusting the application's user interface based on the user's current busy hours (e.g., transitioning to a simplified ``Rush Hour'' mode).
	\item \textbf{Contextual visualization:} Simplifying graphs when users report high stress levels to avoid information overload.
	\item \textbf{Personalized prompts:} Adjusting chatbot questions based on previous mood patterns (e.g., ``Feeling tired after last night's side gig, huh?'').
\end{itemize}

This adaptive approach remains grounded in Nielsen's usability principles (e.g., flexibility and efficiency of use) and iterative design, while adding a contextual layer to enhance effectiveness for users with specific characteristics like non-traditional workers \cite{jmir2025digitalinterventions}. This concept addresses gaps in existing mHealth applications, which often lack adaptability to the cognitive constraints of gig workers, as will be further discussed in the review of prior studies.

\subsection{Supporting Technology: Gemma 3 LLM Architecture}

This study utilizes the Gemma 3 4B model via API as the core conversational engine for Natural Language Understanding (NLU) and response generation. The architecture operates on a dynamic prompt-engineering framework rather than a rigid intent-matching structure.

With its massive 128K context window, the Gemma 3 model is capable of maintaining multi-turn conversational context, retaining user mood history, and generating natural, empathetic, and highly contextual responses. Unlike traditional intent-based chatbots (such as Dialogflow or Rasa) that rely on strict regular expressions or pre-defined conversational flows, the LLM approach allows the system to process unstructured user input flexibly \cite{oulasvirta2024personalization}.

This architecture is highly suitable for providing basic insights based on the user's mood patterns, generating supportive responses such as ``Looking at your trend tonight, it seems you're exhausted after your side gig shift---would you like to talk more about it?'' The Gemma 3 4B API was chosen for its lightweight efficiency for mobile integration, allowing rapid iteration within the iterative design approach while still providing advanced NLU capabilities.

When a user submits an input, the mobile client constructs a payload consisting of three elements: (1) the strict System Prompt (acting as the ethical guardrail and defining the bot's empathetic persona), (2) the user's recently logged mood data (acting as contextual memory), and (3) the user's raw text input. This payload is processed by the Gemma 3 4B model, which leverages its extensive context window to generate a highly personalized, natural language response. This architecture completely replaces modular, regular expression (regex) rule-based systems, allowing the application to parse complex, unstructured human emotion with significantly higher accuracy and conversational fluidity.

\subsection{Psychological Framework of the Chatbot Intervention}

The conversational agent in this study does not aim to provide clinical therapy or psychiatric diagnosis. Instead, the chatbot's interactions are grounded in the psychological principles of Empathetic Active Listening and simplified Cognitive Behavioral Psychoeducation. Active listening in a digital context involves validating the user's reported emotions without immediate judgment or unsolicited problem-solving, which has been shown to significantly reduce acute emotional distress \cite{miner2016chatbot}. Furthermore, by prompting users to identify the specific triggers of their burnout (e.g., ``What made you feel stressed during your shift today?''), the chatbot facilitates self-reflection. This aligns with the self-monitoring principles of Cognitive Behavioral Therapy (CBT), where identifying the relationship between environmental triggers and emotional responses is the first step in effective stress management and emotional regulation for non-traditional workers \cite{beck2020cbt}.

\subsection{Ethical Considerations in Mental Health Applications}

The development and evaluation of mental health applications must adhere to strict ethical principles, including data privacy (encryption and anonymity), informed consent, avoiding clinical diagnoses, and avoiding responses that could be harmful or create dependency \cite{who2023ethics, torous2021standards}. This study applies guidelines from ISO 9241-210 and WHO recommendations for digital mental health interventions by ensuring all mood data is securely stored in Firebase (end-to-end encryption), users are given clear explanations through informed consent, and the chatbot is restricted to supportive assistance without therapeutic claims.

This ethics protocol also includes algorithmic guardrails implemented via System Prompting on the Gemma 3 API to prevent negative or risky responses. If preprocessing filters detect critical trigger words indicating severe psychological distress, the LLM generation is strictly bypassed, and a hardcoded mechanism provides direct referrals to UGM psychological counseling services and emergency hotline numbers. Furthermore, an opt-out mechanism is available at any time, ensuring the research remains safe and responsible for student participants at risk of burnout.

\section{Prior Literature Review}

Table \ref{tab:literature-review} presents a summary of prior studies related to gig economy, burnout, mHealth interventions, and adaptive UI/UX design relevant to this research.

\begin{table}[p]
\centering
\small
\setlength{\tabcolsep}{3pt}
\renewcommand{\arraystretch}{0.95}
\caption{Summary of Prior Studies}
\label{tab:literature-review}
\begin{tabular}{>{\raggedright\arraybackslash}p{0.20\textwidth}
				  >{\centering\arraybackslash}p{0.05\textwidth}
				  >{\raggedright\arraybackslash}p{0.13\textwidth}
				  >{\raggedright\arraybackslash}p{0.17\textwidth}
				  >{\raggedright\arraybackslash}p{0.19\textwidth}
				  >{\raggedright\arraybackslash}p{0.17\textwidth}}
\hline
\textbf{Title} & \textbf{Year} & \textbf{Author(s)} & \textbf{Methods} & \textbf{Main Findings} & \textbf{Gap / Limitations} \\
\hline
Good Gig, Bad Gig: Autonomy and Algorithmic Control in the Gig Economy & 2019 & Wood et al. & Qualitative interviews ($n=150$, global including Indonesia) & Gig work offers flexibility but algorithmic control causes high stress and burnout. & Lacks focus on digital interventions; global sample, not specific to young Indonesians. \\
\hline
World Employment and Social Outlook & 2025 & ILO & Secondary data analysis from global surveys including Indonesia & 41\% of gig workers in Indonesia experience income uncertainty, increasing burnout risk by 50--60\%. & Aggregate data; lacks details on mHealth interventions for gig workers. \\
\hline
Gambaran Burnout pada Mahasiswa Pekerja Paruh Waktu & 2024 & Syiah Kuala Psychology Journal & Cross-sectional survey ($n=200$ part-time students in Aceh) & 91\% of students with side gigs experience burnout due to schedule conflicts; main factor: algorithmic pressure. & No evaluation of digital interventions; limited local sample, not national. \\
\hline
Digital Interventions for Gig Workers & 2025 & JMIR Human Factors & Scoping review of 50+ studies & mHealth apps like mood trackers reduce stress in gig workers by up to 20\%, but lack adaptability. & Minimal studies on adaptive UI; global focus, gap in Indonesian context. \\
\hline
Understanding Communication Patterns in Coping with Stress in Indonesia's Gig Economy & 2025 & Juddi et al.\ (Unpad) & Qualitative phenomenological (interviews $n=5$ gig drivers in Bandung) & Mutual-aid communities reduce stress, but platform algorithms cause exhaustion. & Lacks digital interventions; small sample, suggests multi-case studies with government/platform involvement. \\
\hline
Work Support and Resilience: Mediating Roles of Stress and Health Among Indonesian Gig Workers & 2025 & ResearchGate & SEM-PLS survey ($n=200$ gig workers) & Stress mediates the relationship between support and resilience; gig workers vulnerable to health issues. & Non-digital focus; needs longitudinal studies for tech interventions. \\
\hline
The Influence of Digital Platforms on Gig Workers: A Systematic Literature Review & 2025 & Alauddin et al. & PRISMA review (18 articles, 2019--2023) & Platforms offer flexibility but algorithmic control and lack of protection cause burnout. & Lacks focus on Indonesian region; gap in regulation and adaptive interventions. \\
\hline
What Drives Gig Worker Success? Investigating the Impact of Relevant Experience and Self-Directed Learning & 2025 & Putra et al.\ (UMY) & SEM-PLS survey ($n=235$ freelancers, Indonesia--Malaysia) & Experience and self-learning improve performance, but intrinsic motivation does not moderate. & Cross-sectional; gap in longitudinal studies and digital tools for learning. \\
\hline
Analyzing the Effects of Workload, Work Flexibility, Income Security, and Technological Support on Gig Workers' Productivity & 2025 & Devi et al. & Multiple regression survey ($n=$ gig workers in Jabodetabek) & Technological support dominantly reduces stress and increases productivity. & Not specific to adaptive UI; urban sample, gap in rural areas. \\
\hline
Vulnerable Flexibility: An Analysis of the Representation of Gig Workers' Sustainability in Indonesian Digital Media & 2025 & Murda et al. & Qualitative CDA (120+ media texts, 2023--2025) & Media glorifies flexibility but ignores vulnerabilities; needs cultural recognition. & Focus on representation; gap in empirical mHealth interventions. \\
\hline
The Gig Economy as a Stepping Stone: Career Trajectories of Online Platform Workers in Indonesia & 2023 & Permana (LSE) & Interviews ($n=50$ Indonesian freelancers) & Gig work as a stepping stone but algorithmic control causes isolation and burnout. & Career focus; gap in digital health interventions. \\
\hline
Transformative Dynamics of the Gig Economy: Technological Impacts, Worker Well-Being and Global Research Trends & 2025 & Srihita et al. & Bibliometric analysis (549 Scopus documents) & Technology improves well-being but causes precarity; emerging themes: worker well-being. & Relies on Scopus; gap in sector-specific and algorithmic ethics. \\
\hline
Mood Tracker App With AI Insights: A UX/UI Case Study & 2025 & Reddy & UX/UI case study (design prototype) & Adaptive AI insights increase emotional awareness; playful UI reduces stress. & Not empirical; gap in real-user evaluation for gig workers. \\
\hline
UX for Mental Health: Designing Apps to Promote Wellness and Therapy, Latest Trends in 2025 & 2025 & Vrunik Design & Review of UX trends & Personalization and calming elements are effective; examples include Happify mood tracking. & Non-empirical; gap in Indonesian studies and gig-specific context. \\
\hline
Enhancing Wellbeing Through a UX/UI Optimized Mental Health App & 2025 & Saithong et al. & Pre-post prototype evaluation ($n=50$ students) & Adaptive UX/UI increases resilience and engagement; positive correlation. & Prototype phase; gap in cultural adaptation and long-term effects. \\
\hline
The Influence of Digital Burnout, Screen Time, and Self Control on the Mental Health of Gen Z & 2025 & Nugraheni et al. & SEM-PLS survey ($n=100$ Jakarta students) & Digital burnout and screen time predict poor mental health; stress mediates. & Cross-sectional; limited sample, gap in adaptive interventions. \\
\hline
Effectiveness of Digital Interventions in Reducing Occupational Stress & 2024 & Indra et al.\ (PMC) & Systematic review (15 RCTs) & Digital apps reduce stress for 3--12 months; app-based interventions dominant. & High heterogeneity; gap in LMICs like Indonesia. \\
\hline
A Culturally Adapted Internet-Delivered Mindfulness Intervention & 2025 & Listiyandini et al. & RCT ($n=156$ Indonesian students) & Online mindfulness reduces distress ($g=0.48$--$1.18$); sustained for 1 month. & Short follow-up; gap in cost-effectiveness and adaptive UI. \\
\hline
\end{tabular}
\end{table}

From the review of prior studies, it is evident that gig and non-traditional workers in Indonesia, including students with side gigs or part-time jobs, face a high risk of burnout due to factors such as algorithmic control, income uncertainty, social isolation, and schedule conflicts \cite{wood2019goodgig, juddi2025gigcommunication, syiahkuala2024burnout}. Several studies have demonstrated the effectiveness of digital interventions, such as mHealth applications, mood trackers, and chatbots, in reducing stress symptoms while enhancing self-awareness and resilience \cite{jmir2025digitalinterventions, listiyandini2025mindfulness, saithong2025mhapp}. However, most of these studies remain general, cross-sectional, or focused on global populations, with limited exploration of adaptive UI/UX that specifically addresses cognitive constraints and irregular work schedules among young gig workers in Indonesia.

The primary gap identified is the scarcity of local studies integrating chatbot-based mood trackers with adaptive elements---such as one-tap input, voice options, and contextual personalization---along with evaluations using metrics that encompass emotional and personalization aspects (e.g., ToT, ERI, and PEI). Existing research more frequently assesses general effectiveness without accounting for high cognitive load or the flexible yet unpredictable schedules of non-traditional workers.

Therefore, this study addresses this gap by designing an adaptive UI/UX prototype for a Gemma 3-based mood tracker application specifically targeted at Universitas Gadjah Mada (UGM) students with side gigs, as a subset of young gig workers. The evaluation employs a hybrid approach (in-person observed sessions combined with asynchronous 7-day usage) using three key metrics: ToT for task efficiency under time constraints, ERI for positive emotional impact, and PEI for personalization effectiveness. This is expected to provide novel contributions to digital mental health interventions for non-traditional workers in Indonesia.
